\section{Discussion}
\label{sec:discussion}

\textbf{What remains for devices away from home.} DriftGate gains most where the device exit is strong and gives up some accuracy where the edge exit is strong. A rule that knew the kind of every request, and answered Main requests with the device exit and all other requests with the edge exit, would reach 72.9\% in S1 over the first day, compared with 68.3\% for DriftGate. Most of this gap lies in the OOP and OOR requests of devices away from home. Closing it requires a signal that recognizes, for a single request, that its class is new to the device. The confidence of the device exit gives only a weak signal (Section~\ref{sec:choose}), and the label-free features we combined with a logistic regression did not close the gap in our runs. Distances in the feature space of the device block, measured against the device's own training data, are a natural next candidate.

\textbf{Trained models.} The first-day results of Table~\ref{tab:main} include the early phase of training, when the edge exit is weak and the device exit carries more of the accuracy. With trained models (Section~\ref{sec:trained}), the gap between the two exits and the other combining rules narrows, and the parts of DriftGate that help depend on the scenario. In the S1 replay, the corrected edge exit alone and a fixed weight of 0.2 are more accurate than DriftGate, and so is the uncorrected probability average. In the S2 replay, DriftGate remains the most accurate rule. A deployment that serves trained models over many days should therefore check the combination and the weight on its own traffic, for example on a development day. Our replay reuses the same day of movement for S2's real traces, so it is not an evaluation on a new day.

\textbf{Relation to training.} DriftGate acts after training and only on the answers of the requesting device, and it does not change the shared model. The location-aware training policy of Section~\ref{sec:training_side} showed a gain for devices at home together with a loss for devices away, and we have not tested whether combining it with DriftGate keeps both effects. Because DriftGate needs only the two output vectors, the class list of the device, and the class counts of the cell, it can run on top of other ways of training the two exits, such as different personalization ratios, personalized federated learning objectives, or the online distillation of Kim \emph{et al.}~\cite{kim2025onlinepsfl}; we have not measured its accuracy with these.

\textbf{Information at the edge.} The prior correction uses the device's class list $M_k$ and the class counts of the cell. In split learning, the edge receives training labels and already holds both. In variants that keep labels on the device, the device can send its class list once, which reveals no more than which classes it trains on.

\textbf{Scope.} We studied changes in which classes a device sees, which is how mobility changes the traffic of a recognition application in our simulation. Changes in how the same classes look, such as lighting or sensor differences between places, affect both exits and are outside the scope of the correction. The class mix of the requests follows a simulation setting ($\rho$), and the real GeoLife traces determine only where devices are.
